% ======================================================================
%  KAVRAM: MCP — Model Context Protocol  (Tier 1)  — kompakt
% ======================================================================

\concept{Actions · Eylem}{MCP}{Modeli araçlara bağlamanın ortak standardı: AI için USB-C}

\stage{1}{Problem}
Her araç, her uygulamaya elle bağlanıyor. 3 farklı AI uygulaman ve 10 aracın
varsa, $3\times10$ özel entegrasyon yazıyorsun --- her biri ayrı bakım, ayrı
hata kaynağı. Bir aracı bir yerde bağlaman, onu diğer uygulamada işe yaramaz
kılıyor. Standart yok.

\keyline{Sorun araçlarda değil, her bağlantının \emph{tekil ve tekrar kullanılamaz} olmasında. $N\times M$ entegrasyon sürdürülemez.}

\stage{2}{Neden Var}
MCP, modellerin araçlara, verilere ve hazır prompt'lara nasıl bağlanacağını
\emph{standartlaştırır}. Bir entegrasyonu bir kez ``MCP sunucusu'' olarak yaz;
MCP uyumlu her istemci (IDE, sohbet uygulaması, agent) onu yeniden yazmadan
kullansın. $N\times M$ problemi $N+M$'e iner.

\stage{3}{Tarihçe}
\begin{timeline}
\yr{2023} özel function calling: her uygulama kendi entegrasyonu \tlsep \yr{2023} kapalı ``plugin'' ekosistemleri \tlsep \yr{2024} MCP: açık protokol olarak duyurulur \tlsep \yr{2024–26} sunucu ekosistemi büyür (dosya, DB, SaaS araçları)
\end{timeline}

\stage{4}{Tanım}
\begin{problembox}
\textbf{MCP (Model Context Protocol):} AI uygulamalarının araçları, kaynakları
ve prompt'ları nasıl \emph{sunup tükettiğini} standartlaştıran açık protokol.
\textbf{Sunucu} (server) yetenekleri sunar; \textbf{istemci} (client, AI
uygulaması) onları tüketir --- ortak bir dille.
\end{problembox}

\stage{5}{Zihinsel Model}
\begin{mentalbox}
\textbf{USB-C} gibi düşün. Eskiden her cihazın kendi kablosu vardı --- telefon,
kamera, yazıcı; hepsi farklı. USB-C tek standart bağlantı noktası getirdi. MCP,
AI uygulamaları ile ihtiyaç duydukları araç/veri arasındaki o \emph{standart
porttur}: bir kez tak, her yerde çalışsın.
\end{mentalbox}

\stage{6}{Mimari Diyagram}
Tek protokol, çok istemci ile çok sunucuyu birbirine bağlar.
\vspace{3pt}
\begin{center}
\begin{tikzpicture}[font=\sffamily\scriptsize, node distance=4mm,
   cl/.style={draw=accentdk, fill=accent!12, rounded corners=2pt, align=center,
              minimum height=8mm, minimum width=24mm},
   sv/.style={draw=accent, fill=wash, rounded corners=2pt, align=center,
              minimum height=8mm, minimum width=24mm},
   bus/.style={draw=accentdk, fill=accentdk, text=white, rounded corners=2pt,
               align=center, minimum height=8mm, minimum width=70mm},
   flow/.style={-{Stealth[length=1.8mm]}, accent, line width=0.7pt}]
  % istemciler
  \node[cl] (c1) {IDE};
  \node[cl, right=5mm of c1] (c2) {Sohbet uygulaması};
  \node[cl, right=5mm of c2] (c3) {Agent};
  % protokol barı
  \node[bus, below=7mm of c2] (bus) {MCP \,(ortak protokol)};
  % sunucular
  \node[sv, below=7mm of bus.south west, anchor=north west] (s1) {Dosya sistemi};
  \node[sv, right=5mm of s1] (s2) {Veritabanı};
  \node[sv, right=5mm of s2] (s3) {SaaS araçları};

  \foreach \c in {c1,c2,c3} \draw[flow] (\c) -- (bus);
  \foreach \s in {s1,s2,s3} \draw[flow] (bus) -- (\s);
  \node[font=\sffamily\tiny, text=softgray, left=1mm of c1] {istemciler};
  \node[font=\sffamily\tiny, text=softgray, left=1mm of s1] {sunucular};
\end{tikzpicture}
\end{center}
\keyline{Her istemci her sunucuyla otomatik konuşur: $N\times M$ özel entegrasyon yerine $N+M$ standart bağlantı.}

\stage{7}{Karar Çerçevesi}
MCP mi, doğrudan function calling mi? Tekrar kullanım ve ölçek belirleyici.\par
\vspace{2pt}
\noindent\begin{tabularx}{\linewidth}{@{}>{\raggedright\arraybackslash}p{40mm} L@{}}
\arrayrulecolor{hairline}\toprule
{\sffamily\bfseries\small Durum} & {\sffamily\bfseries\small Seçim} \\
\midrule
Araç birçok uygulamada kullanılacak & \textbf{MCP sunucusu} \\
Çok sayıda entegrasyon / ekosistem & \textbf{MCP} \\
Tek uygulamada tek, basit araç & Doğrudan function calling \\
Hızlı prototip / tek seferlik & Doğrudan function calling \\
\bottomrule
\end{tabularx}

\stage{8}{Üretim Senaryosu}
\begin{prodbox}[Üretim Senaryosu · Paylaşılan İç Sistemler]
\textbf{Bağlam.} Bir şirkette CRM, dosya deposu ve analitik veritabanı var. Üç
farklı AI uygulaması (IDE asistanı, destek botu, iç agent) hepsine erişmek
istiyor --- ama her biri ayrı ayrı entegre ediliyor, bakım kabusu.

\smallskip
\textbf{Çözüm.} Her iç sistem bir kez \textbf{MCP sunucusu} olarak sunulur.
Üç uygulama da MCP istemcisi olarak aynı sunucuları yeniden entegrasyon olmadan
kullanır.

\smallskip
\textbf{Sonuç.} Yeni bir AI uygulaması eklemek ``sıfırdan entegrasyon'' değil,
``mevcut sunuculara bağlan'' demek. Bakım ve güvenlik tek noktada toplanır.
\end{prodbox}

\stage{9}{Anti-pattern}
\begin{antibox}[Anti-pattern · ``Her Şeye MCP / Açık Sunucu'']
Tek, basit bir aracı koca bir MCP sunucusuna dönüştürerek gereksiz karmaşa
yaratmak; ya da daha tehlikelisi, MCP sunucusunu yetkilendirme ve sandboxing
olmadan açmak --- böylece modele dosya sistemine veya veritabanına sınırsız
erişim vermek.

\smallskip
\textbf{Belirti:} Basit iş için aşırı altyapı; veya bir sunucu üzerinden geniş, kontrolsüz erişim.\\
\textbf{Kök neden:} Standardı amaç sanmak; güvenliği protokolden beklemek.\\
\textbf{Düzeltme:} MCP'yi tekrar kullanım varken kullan. Yetki ve sınırı her zaman \emph{sunucunun içinde} uygula, sandbox'la.
\end{antibox}

\stage{10}{Alternatifler \& Ödünleşimler}
\vspace{2pt}
\begin{tabularx}{\linewidth}{@{}>{\raggedright\arraybackslash}p{30mm} L L@{}}
\arrayrulecolor{hairline}\toprule
{\sffamily\bfseries\small Yaklaşım} & {\sffamily\bfseries\small Güçlü} & {\sffamily\bfseries\small Ödünleşim} \\
\midrule
\textbf{MCP} & Tekrar kullanım, ekosistem, standart & Kurulum yükü; olgunlaşan protokol \\
\addlinespace[2pt]
\textbf{Doğrudan function calling} & Basit, hızlı, tek uygulama & Her entegrasyon tekrar yazılır \\
\addlinespace[2pt]
\textbf{Özel plugin sistemi} & Tam kontrol & Kapalı; taşınmaz \\
\addlinespace[2pt]
\textbf{Düz API entegrasyonu} & Her yerde çalışır & Model-farkında değil; tutkal kod \\
\bottomrule
\end{tabularx}

\stage{11}{İlişkili Kavramlar}
\term{Tool Calling} (MCP'nin standartlaştırdığı şey) \tlsep \term{A2A} (agent'lar arası protokol) \tlsep
\term{API} (alttaki arabirim) \tlsep \term{Agent} (MCP sunucularını tüketen) \tlsep
\term{Authorization} (sunucu içinde yetki)

\stage{12}{Özet}
\begin{itemize}
  \item MCP, model--araç bağlantısını standartlaştırır: AI için USB-C.
  \item $N\times M$ özel entegrasyonu $N+M$ standart bağlantıya indirir.
  \item Sunucu yetenek sunar, istemci tüketir --- ortak dille.
  \item Tekrar kullanım/ölçek varsa MCP; tek seferlik basit araçta doğrudan function calling.
  \item Güvenlik protokolden gelmez: yetki ve sandbox sunucunun içindedir.
\end{itemize}

\stage{13}{Mühendislik İlkesi}
\begin{principlebox}
\textbf{İlke \#8 — Entegrasyonu bir kez yaz, her yerde kullan; ama sınırı sunucuda uygula.}\\[4pt]
\normalsize Standart, tekrarı ortadan kaldırır --- güvenliği değil. Her MCP
sunucusu kendi yetki ve sandbox sınırını taşımalıdır.
\end{principlebox}
